\documentclass[a4paper]{article}
\usepackage[pdftex]{graphicx}
\usepackage[utf8]{inputenc}
\usepackage{enumerate}
\usepackage{amssymb}
\usepackage{icomma}
\usepackage{siunitx}
\sisetup{locale=DE} 
\usepackage{href-ul}
\hypersetup{
	colorlinks=true,
	linkcolor=blue,
	urlcolor=blue}
\usepackage{geometry}
\geometry{a4paper, top=15mm, left=15mm, right=15mm, bottom=15mm,
	headsep=10mm, footskip=12mm}

\begin{document}
{\bf Druck, Kraft und Arbeit}


{\bf Vorübung: }

\renewcommand{\arraystretch}{1.5}
\begin{tabular}{|p{1.4 cm}|p{1.4 cm}|p{1.4 cm}|p{1.4 cm}|p{1.4 cm}|}
	\hline
	& $\SI{}{\meter}$ & $\SI{}{\deci\meter}$ & $\SI{}{\centi\meter}$ & $\SI{}{\milli\meter}$\\
	\hline
	$\SI{1}{\meter}$ & 1 & & &\\
	\hline
	$\SI{1}{\deci\meter}$ & & 1 & &\\
	\hline
	$\SI{1}{\centi\meter}$ & & & 1 &\\
	\hline
	$\SI{1}{\milli\meter}$ & & &  & 1\\
	\hline
	
\end{tabular}
\hspace{0.5 cm}
\begin{tabular}{|p{1.4 cm}|p{1.4 cm}|p{1.4 cm}|p{1.4 cm}|}
	\hline
	& $\SI{}{\pascal}=\SI{}{\newton\over\meter^2}$ & $\SI{}{\kilo\pascal}$ & bar\\
	\hline
	$\SI{1}{\pascal}$ & 1 &  & \\
	\hline
	$\SI{1}{\kilo\pascal}$ &  & 1 & \\
	\hline
	1 bar & & & 1 \\
	\hline
\end{tabular}

\vspace{0.5 cm}
\begin{tabular}{|p{1.8 cm}|p{1.8 cm}|p{1.8 cm}|p{1.8 cm}|}
	\hline
	& $\SI{}{\meter^2}$ & $\SI{}{\deci\meter^2}$ & $\SI{}{\centi\meter^2}$ \\
	\hline
	$\SI{1}{\meter^2}$ & 1 & & \\
	\hline
	$\SI{1}{\deci\meter^2}$ & & 1 & \\
	\hline
	$\SI{1}{\centi\meter^2}$ & & & 1\\
	\hline
\end{tabular}
\hspace{0.5 cm}
\begin{tabular}{|p{1.8 cm}|p{1.8 cm}|p{1.8 cm}|}
	\hline
	&$\SI{}{\newton}$ & $\SI{}{\kilo\newton}$  \\
	\hline
	$\SI{1}{\newton}$ & 1 & \\
	\hline
	$\SI{1}{\kilo\newton}$ & & 1	\\
	\hline
\end{tabular}


\begin{minipage}{0.3\textwidth}
	Das Bild zeigt die Wirkungsweise einer Hydraulik.
	Die Arbeit $$W=F\cdot s$$ ist auf beiden Seiten gleich,
	ebenso ist der Druck $$p={F \over A}$$  überall gleich.
	Weiter gilt:
	$${F_2 \over F_1}={A_2 \over A_1}={s_1 \over s_2}$$
\end{minipage}
\hfill
\begin{minipage}{0.7\textwidth}
	\includegraphics[width=11 cm]{hydr076}
\end{minipage}


\begin{enumerate}[1.]

{\bf \item Fülle die Tabelle aus!}

\renewcommand{\arraystretch}{2}
\begin{tabular}{|p{1.8 cm}|p{1.8 cm}|p{1.8 cm}|p{1.8 cm}|p{1.8 cm}|p{1.8 cm}|p{1.8 cm}|p{1.8 cm}|}
\hline
Kraft 1 & Fläche 1 & Weg 1 & Kraft 2 & Fläche 2 & Weg 2 & Druck & Arbeit \\ 
$F_1$ & $A_1$ & $s_1$ & $F_2$ & $A_2$ & $s_2$ & $p$ & $W$ \\
\hline
$\SI{12}{\newton}$ & $\SI{8}{\centi\meter^2}$ &$\SI{0,70}{\meter}$ &$\SI{2,4}{\kilo\newton}$ & & & & \\
\hline
$\SI{80}{\newton}$ & $\SI{30}{\centi\meter^2}$ & & &$\SI{0,24}{\meter^2}$ &  & &$\SI{48}{\joule}$ \\
\hline
 &  &$\SI{0,15}{\meter}$ &$\SI{7,0}{\newton}$ & &$\SI{3,0}{\centi\meter}$ &$\SI{7,0}{\kilo\pascal}$ & \\
\hline
$\SI{0,20}{\kilo\newton}$ &  &$\SI{0,40}{\meter}$ & &$\SI{0,20}{\meter^2}$ &$\SI{8,0}{\milli\meter}$ & &\\
\hline
 & $\SI{4,0}{\centi\meter^2}$ & & & &$\SI{34}{\centi\meter}$ &$\SI{13}{\kilo\pascal}$ &$\SI{8,8}{\joule}$ \\
\hline
\end{tabular}



\fbox{
	\begin{minipage}{0.5\textwidth}
		Zur Lösung bitte \href{https://www.okuyakl.de/phys/p9drkL076/ll076.pdf}{hier klicken} oder den QR-Code scannen.\\
	Weitere Arbeitsblätter gibt es unter 
	
	\href{https://www.okuyakl.de}{www.okuyakl.de}
	\end{minipage}
	\hfill
	\begin{minipage}{0.4\textwidth}
		\includegraphics[width=1.5 cm]{../../viecher/zwe03}
		\includegraphics[width=3 cm]{qr076}
		\includegraphics[width=2 cm]{../../viecher/afanticon1}
		
\end{minipage}}



\end{enumerate} 

\end{document}

